\documentclass{article}

\usepackage[utf8]{inputenc}
\usepackage[T2A]{fontenc}
\usepackage[english, russian]{babel}
\usepackage{amsmath}
\usepackage{amsfonts}
\usepackage{mathtools}
\usepackage{datetime}
\usepackage[pdftex]{graphicx}
\usepackage{indentfirst}
\usepackage{asymptote}
\usepackage{systeme}
\usepackage{hyperref}
\hypersetup{
	colorlinks=true,
	linkcolor=blue,
	filecolor=magenta,      
	urlcolor=cyan,
	pdftitle={Overleaf Example},
	pdfpagemode=FullScreen,
}

\title{\bf Лабораторная часть \\ Экстремумы ФНП}
\author{Мовчан Игорь Евгеньеви}
\newdate{date}{19}{11}{2023}
\date{\displaydate{date}}
\graphicspath{{./}}

\begin{document}
	\maketitle
	\section{Аналитический метод (продолжение)}
	Функция задана следующим соотношением (8 вариант заданяи)
	$$
	z(x, y) = x^2 + y^2 -x^3 - y^3 + 2xy.
	$$
	
	Её график вместе с найденными стационарными точками $(\frac{4}{3}, \frac{4}{3})$ и $(0, 0)$, выглядит, соответственно, следующим образом:
	\begin{center}
		\includegraphics[scale=0.48]{alo.png}
	\end{center}
	
	Визуализация неравенства $x^3 + y^3 > x^2 + y^2 + 2xy$ представлена ниже:
	\begin{center}
	\includegraphics[scale=0.35]{alo1.png}
	\end{center}

	
	В пространстве оно полностью покрывает ту часть графика, которая находится ниже z = 0, чего и хотелось достичь:
	\begin{center}
		\includegraphics[scale=0.5]{alo2.png}
	\end{center}

	\section{Численный метод}
	Выберем начальную (стартовую) точку $(5, 5) = (x_0, y_0)$, на каждом шаге будем вычислять значение функции, приращение функции и градиент в текущей точке $(x_k, y_k)$, определять следующую точку по формуле $(x_{k+1}, y_{k+1}) = (x_k, y_k) + a_k \text{grad} f (x_k, y_k)$. Суммируем, так как градиент показывает наискорейший рост функции, а мы как раз-таки ищем локальный максимум функции.
	
	Градиент функции будет задаваться вектором 
	$$
	\text{grad} f (x, y) =
	\begin{bmatrix}
		2x - 3x^2 + 2y \\
		2y - 3y^2 + 2x
	\end{bmatrix}
	$$
	
	Посмотрим теперь на графики. Во всех ситуациях чем точка зеленее, тем она дальше ушла по итерации, а чем краснее, тем, соответственно, ближе к началу вычислений. Также на графиках будет прописываться критерий остановки, если <<функция>> - по малости приращений значения функции, то есть $|\Delta f|<\varepsilon = 0.001$, а если <<аргумент>>, значит, по малости приращения аргумента $||(\Delta x, \Delta y) || < d = 0.001$. Рассматриваемая норма для расчёта приращения аргумента была стандартной теоремой Пифагора.
	
	Рассмотрим ситуацию $a_k = 0.05$ (постоянное, малое значение, последнее нужно, чтобы значения не улетали на небеса):
	\begin{center}
		\includegraphics[scale=0.43]{Figure_1.png}
		\includegraphics[scale=0.43]{Figure_2.png}
	\end{center}

	Видим, что значения достаточно близки к найденным вручную, а значит, программа хотя бы работает. Также графики показывают, что вычисления через функцию и аргумент сильно близки и при достойном $a_k$ оба дают хорошие приближения.
	
	Поэкспериментируем с $a_k$, рассмотрим <<линейное убывание>>. Достаточно приближенные к исходным точки значения получаем при $a_k = \frac{1}{20}(1-\frac{k}{100})$:
	\begin{center}
		\includegraphics[scale=0.43]{Figure_3.png}
		\includegraphics[scale=0.43]{Figure_4.png}
	\end{center}
	
	\break
	В ином же случае получаем либо взлёт по значениям, а это означает, что шаг выбран слишком большим, либо непохожие на реальность результаты (ниже представлены графики для $a_k = \frac{1}{200}(1-\frac{k}{100})$):
	\begin{center}
		\includegraphics[scale=0.43]{Figure_5.png}
		\includegraphics[scale=0.43]{Figure_6.png}
	\end{center}
	
	Это говорит нам о том, что $a_k$ не должны быть как слишком малыми, так и слишком большим. И в том, и в том случае результаты работы программы нас будут неудовлетворять. Стоит отметить, что важным критерием подбора малости было количество итераций. Их малое количество свидетельствует о слишком малом $a_k$, то есть условие остановки срабатывает слишком рано. В идеале (исходя из опыта) их должно быть не меньше 15. Их также не должно быть слишком много (как мы увидели из примера с линейной $a_k$), ведь это тоже даёт нам неудовлетворительный результат.
	
	\break
	Наконец, изучим ситуацию, когда $a_k$ представлена в виде убывающей геометрической прогрессии. Возьмём $a_k = \frac{1}{15}(\frac{9}{10})^k$, где $\frac{1}{15}$ можно назвать нормировочный коэффициентом, нужным для того для того, чтобы функция не улетала на большие значения в самом начале, когда показатель в степени $k$ ещё не слишком большой. В итоге получаем следующие графики:
	\begin{center}
		\includegraphics[scale=0.43]{Figure_7.png}
		\includegraphics[scale=0.43]{Figure_8.png}
	\end{center}
	
	Полученное полностью согласуется с полученными в аналитической части результатами.
	
	Код программы вместе со всеми графиками: \href{https://github.com/ventusjke/lab_gradient}{github link}
	
	
	
	
	
	
	
	
	
	
		
\end{document}
